\subsection{Reentrancy}
\label{sec:reentrancy}

Move on Aptos later gained \emph{function values}: a closure that can be captured, passed around, stored (into the global storage), and invoked.
This feature greatly enhanced the expressivity of the language (by enabling dynamic dispatch), but also reintroduced a classic challenge -- \emph{reentrancy}.

Prior to this, freedom from reentrancy was a \emph{static} guarantee: a function declares its global resource accesses with the |acquires| annotations, 
and because module dependencies must be acyclic, no function a module calls can re-enter that module and borrow the same resource twice.

Here is an example demonstrating what a vulnerable pattern may look like, using a withdrawal function that processes the payout through a caller-supplied
callback:

\paragraph{Example.}
\begin{MoveBox}
module 0x42::vault {
  struct Vault has key { balance: u64 }

  public fun withdraw(owner: address, payout: |u64|) acquires Vault {
    let amount = Vault[owner].balance;
    payout(amount);            // dynamic dispatch into unknown code
    Vault[owner].balance = 0;  // cleared only after the callback
  }
}
\end{MoveBox}

A malicious or poorly written |payout| could call |withdraw| again before the first withdrawal has been completed, resulting in the re-entrant call reading the same non-zero |balance|.

\paragraph{Runtime Reentrancy Check.}
Our current solution to this problem is a runtime reentrancy check. The VM's runtime maintains a map that for each module, records a count of how many times it has been entered.
The count is incremented when entered from a different module, or from a function value, even if the underlying function is in the same module, and decremented on return.

A count greater than one is not in itself a violation. A transaction-aborting error is only fired when 
(1) a global resource defined by the reentered module is accessed, or 
(2) a module is reentered while previously locked by a \emph{module lock}, which is a stricter mode activated by entering a function with the |#[module_lock]| attribute.

Unlike the type, ability and reference checks, which are redundant safety nets that reassert key properties guaranteed by the bytecode verifier, the reentrancy check is a \emph{required} check,
crucial for ensuring safety with dynamic dispatch. Its implementation however, shares the same spirit as the other runtime checks.

Because the reentrancy check is invoked on every function entry and return, it sits on the interpreter's hot path.
The cost is negligible today thanks to two optimizations (fast, non-cryptographic hash + module id interning), but could be further under test with our next-gen VM, 
which will have much faster function calls, making the per-call overhead more noticeable.

We are also exploring whether the check could instead be deferred to the asynchronous phase (Section~\ref{sec:async}). However its soundness requires more study
as aborts caused by it are part of a transaction's execution semantics, rather than a redundant post-execution validation.

%%% Local Variables:
%%% mode: latex
%%% TeX-master: "main"
%%% End:
